\documentclass[12pt]{article}
\usepackage[utf8]{inputenc}
\usepackage[T1]{fontenc}
\usepackage[margin=1in]{geometry}
\usepackage{setspace}
\usepackage{xcolor}
\usepackage[colorlinks=true, urlcolor=blue!60!black, linkcolor=blue!60!black]{hyperref}
\setlength{\parindent}{2em}
\setlength{\parskip}{1.25em}
\renewcommand{\baselinestretch}{1.0}
\begin{document}
\begin{center}
{\LARGE \textit{Nixon v. United States}, 506 U.S. 224 (1993)\par}
\vspace{0.8em}
{\large Case Brief\par}
\end{center}
\vspace{1em}

\section{Memory Jogger}

Because Article I, \S~3, cl.\ 6 grants the Senate the ``sole Power to try all Impeachments'' and the word ``try'' supplies no judicially manageable standard, a federal judge's constitutional challenge to the Senate's use of a factfinding committee under Rule XI presents a nonjusticiable political question (\href{https://www.courtlistener.com/opinion/112803/nixon-v-united-states/}{CourtListener}; \href{https://supreme.justia.com/cases/federal/us/506/224/}{Justia}).

\section{Detailed Case Facts}

Walter L. Nixon, Jr., served as Chief Judge of the United States District Court for the Southern District of Mississippi (\href{https://www.law.cornell.edu/supct/html/91-740.ZO.html}{Cornell LII}). A federal grand jury investigated reports that Nixon had accepted a gratuity from a Mississippi businessman in exchange for asking a local district attorney to halt the prosecution of the businessman's son. Following that investigation, a jury convicted Nixon of two counts of making false statements before a federal grand jury, and he was sentenced to prison (\textit{United States v. Nixon}, 816 F.2d 1022 (5th Cir. 1987); \href{https://supreme.justia.com/cases/federal/us/506/224/}{506 U.S. at 226}). Because Nixon refused to resign his Article III commission, he continued collecting his judicial salary while incarcerated (\href{https://www.law.cornell.edu/supct/html/91-740.ZO.html}{Cornell LII}).

On May 10, 1989, the U.S. House of Representatives adopted three articles of impeachment charging Nixon with high crimes and misdemeanors: two articles for giving false testimony before the grand jury and a third for bringing disrepute upon the federal judiciary (\href{https://supreme.justia.com/cases/federal/us/506/224/}{506 U.S. at 226--27}). After the House presented the articles to the Senate, the Senate invoked Senate Impeachment Rule XI, which authorizes the Presiding Officer, if the Senate so orders, to appoint a committee of Senators to ``receive evidence and take testimony'' and report a certified transcript to the full Senate (\href{https://www.law.cornell.edu/supct/html/91-740.ZO.html}{Cornell LII}).

The appointed Senate committee conducted four days of evidentiary hearings at which ten witnesses, including Nixon himself, testified (\href{https://supreme.justia.com/cases/federal/us/506/224/}{506 U.S. at 227}). The committee then transmitted to the full Senate a complete transcript of the hearings and a report outlining uncontested facts and summarizing the evidence on contested facts. Both Nixon and the House impeachment managers filed extensive final briefs and presented oral arguments on the Senate floor during three hours allotted for argument on November 1, 1989; Nixon made a personal appeal, and several Senators questioned both sides directly. On November 3, 1989, the full Senate voted by more than the constitutionally required two-thirds majority to convict Nixon on the first two articles, and the presiding officer entered judgment removing him from office (\href{https://supreme.justia.com/cases/federal/us/506/224/}{506 U.S. at 227--28}).

\section{Procedural History}

Following his removal, Nixon filed suit in the United States District Court for the District of Columbia seeking a declaratory judgment that his impeachment conviction was void and requesting reinstatement of his judicial salary and privileges (\href{https://www.law.cornell.edu/supct/html/91-740.ZS.html}{Cornell LII Syllabus}). He argued that Senate Rule XI violated the Impeachment Trial Clause, U.S. Const.\ art.\ I, \S~3, cl.\ 6, because delegating evidentiary hearings to a committee prevented the full Senate from ``trying'' his impeachment.

The District Court dismissed the action as a nonjusticiable political question, 744 F. Supp. 9 (D.D.C. 1990) (\href{https://supreme.justia.com/cases/federal/us/506/224/}{506 U.S. at 228}). The United States Court of Appeals for the District of Columbia Circuit affirmed, 290 U.S. App. D.C. 420, 938 F.2d 239 (1991). The Supreme Court granted certiorari, 502 U.S. 1090 (1992), heard oral argument on October 14, 1992, and affirmed on January 13, 1993 (\href{https://www.oyez.org/cases/1992/91-740}{Oyez}).

\section{Judicial Votes}

The Supreme Court affirmed the judgment of the D.C. Circuit by a unanimous vote of \textbf{9--0}, though the Justices divided across four separate opinions regarding the rationale (\href{https://www.law.cornell.edu/supct/html/91-740.ZS.html}{Cornell LII Syllabus}; \href{https://supreme.justia.com/cases/federal/us/506/224/}{Justia}):

\begin{itemize}
\item \textbf{Majority Opinion (6 Justices):} Chief Justice William H. Rehnquist delivered the opinion of the Court, joined by Justices John Paul Stevens, Sandra Day O'Connor, Antonin Scalia, Anthony M. Kennedy, and Clarence Thomas.
\item \textbf{Concurring Opinion (1 Justice):} Justice John Paul Stevens filed a brief concurring opinion (while also joining the majority opinion in full).
\item \textbf{Concurring in the Judgment (2 Justices):} Justice Byron R. White filed an opinion concurring in the judgment, joined by Justice Harry A. Blackmun, concluding that the claim was justiciable but failed on the merits.
\item \textbf{Concurring in the Judgment (1 Justice):} Justice David H. Souter filed an opinion concurring in the judgment on prudential political-question grounds while leaving open judicial intervention in extreme cases.
\item \textbf{Dissents:} None (0).
\end{itemize}

\section{Holding}

\textbf{Judgment:} Affirmed. The D.C. Circuit's dismissal of Nixon's suit was upheld (\href{https://www.courtlistener.com/opinion/112803/nixon-v-united-states/}{CourtListener}).

\textbf{Rule Announced:} A constitutional challenge to the Senate's procedures for trying an impeachment under Article I, \S~3, cl.\ 6---specifically its use of a committee to take testimony and receive evidence under Senate Rule XI---presents a nonjusticiable political question because the Impeachment Trial Clause contains a textually demonstrable constitutional commitment of the issue to the Senate and the word ``try'' lacks sufficient precision to yield judicially discoverable and manageable standards (\href{https://supreme.justia.com/cases/federal/us/506/224/}{506 U.S. at 228--38}).

\textbf{Holding vs.\ Dicta:} The strict holding is that courts may not adjudicate whether Senate Rule XI violates the Impeachment Trial Clause's directive that the Senate ``try'' impeachments. Chief Justice Rehnquist's broader remarks suggesting that judicial review of impeachment proceedings would undermine finality in the event of a presidential impeachment, create unmanageable remedial difficulties (such as reinstating a judge whose seat has been filled), or involve any judicial role beyond the three express textual requirements of Clause 6 (oath or affirmation, two-thirds concurrence, and the Chief Justice presiding at a presidential trial) sweep beyond the facts of a judicial impeachment under Rule XI. Likewise, the Court did not squarely resolve whether a violation of one of those three explicit textual requirements---or the extreme hypotheticals raised by Justices White and Souter (such as conviction by coin toss)---would be justiciable, though the majority's emphasis on ``identifiable textual limits'' in distinguishing \textit{Powell v. McCormack}, 395 U.S. 486 (1969), implies in dicta that breaching those three explicit limits could remain reviewable (\href{https://supreme.justia.com/cases/federal/us/506/224/}{506 U.S. at 236--38}).

\section{Analysis of Opinions}

\subsection{Majority Opinion (Chief Justice Rehnquist)}

Chief Justice Rehnquist anchored the Court's analysis in the first two prongs of \textit{Baker v. Carr}, 369 U.S. 186, 217 (1962): ``a textually demonstrable constitutional commitment of the issue to a coordinate political department'' and ``a lack of judicially discoverable and manageable standards for resolving it'' (\href{https://supreme.justia.com/cases/federal/us/506/224/}{506 U.S. at 228}). Rehnquist observed that these two criteria are interrelated, because the absence of manageable standards reinforces the inference of a textual commitment.

Turning to the text of Article I, \S~3, cl.\ 6, Rehnquist rejected Nixon's contention that the verb ``try'' mandates a full-body, judicial-style evidentiary hearing. Citing eighteenth-century and modern dictionaries---including Samuel Johnson's 1785 dictionary and \textit{Webster's Third New International Dictionary}---he concluded that ``try'' carries multiple broad meanings and fails to establish an implied procedural mandate (\href{https://supreme.justia.com/cases/federal/us/506/224/}{506 U.S. at 229--30}). This reading was reinforced by structural contrast: the second and third sentences of Clause 6 impose three highly specific requirements (oath or affirmation, a two-thirds vote to convict, and the Chief Justice presiding when the President is tried), indicating that the Framers deliberately omitted additional procedural restrictions.

Rehnquist next gave decisive weight to the word ``sole,'' noting that it appears only twice in the Constitution---both in the Impeachment Clauses (Article I, \S~2, cl.\ 5 and \S~3, cl.\ 6) (\href{https://www.law.cornell.edu/supct/html/91-740.ZO.html}{Cornell LII}). Relying on dictionary definitions of ``sole'' as ``functioning \dots{} independently and without assistance or interference,'' he dismissed Nixon's claim that the Committee of Style's addition of ``sole'' was a mere cosmetic edit, emphasizing that the Convention adopted the final enacted text. He also rejected Nixon's reliance on the Pardon Clause exception (``except in Cases of Impeachment,'' Art.\ II, \S~2, cl.\ 1), explaining that a pardon mitigates criminal punishment rather than overturning a tribunal's judgment.

Finally, Rehnquist invoked constitutional history, separation of powers, and pragmatic concerns. The Constitutional Convention explicitly rejected proposals to place impeachment trials in the Supreme Court, and Alexander Hamilton's \textit{Federalist Nos.\ 65, 66, 79, and 81} confirmed that impeachment was designed as the Legislature's sole constitutional check on the Judiciary (\href{https://supreme.justia.com/cases/federal/us/506/224/}{506 U.S. at 233--36}). Permitting courts to review impeachments would allow the Judiciary to police the very mechanism meant to check it, risk bias in separate criminal prosecutions under Article I, \S~3, cl.\ 7, expose the nation---especially in a presidential impeachment---to prolonged chaos and lack of finality, and raise intractable remedial problems. Rehnquist distinguished \textit{Powell v. McCormack}, 395 U.S. 486 (1969), on the ground that Article I, \S~2 supplied fixed, separate textual qualifications for House membership, whereas no separate constitutional provision defines ``try'' (\href{https://supreme.justia.com/cases/federal/us/506/224/}{506 U.S. at 236--38}).

\subsection{Concurring Opinion (Justice Stevens)}

Justice Stevens joined the majority opinion in full and added a concise concurrence emphasizing institutional trust and judicial restraint (\href{https://www.law.cornell.edu/supct/html/91-740.ZC.html}{Cornell LII Stevens Concurrence}). Rather than resting heavily on fine textual inferences from ``sole'' and ``try,'' Stevens stressed the Framers' overarching choice to entrust impeachment to the Legislative Branch. Citing the Senate's acquittal of Justice Samuel Chase in 1805, he argued that history demonstrates the Senate's solemn awareness of its constitutional responsibility and that respect for a coordinate branch forecloses assuming the Senate would engage in the ``improbable hypotheticals'' raised by Justices White and Souter (\href{https://supreme.justia.com/cases/federal/us/506/224/}{506 U.S. at 238}).

\subsection{Opinion Concurring in the Judgment (Justice White, joined by Justice Blackmun)}

Justice White disagreed sharply with the Court's nonjusticiability holding, arguing that courts have both the authority and the duty to decide whether the Senate complied with the constitutional command to ``try'' impeachments (\href{https://www.law.cornell.edu/supct/html/91-740.ZC1.html}{Cornell LII White Concurrence}). White pointed to the Solicitor General's concession at oral argument that, under the Government's theory, the Senate could constitutionally convict an impeached official without any hearing at all simply by declaring him ``a bad guy'' (\href{https://supreme.justia.com/cases/federal/us/506/224/}{506 U.S. at 239}).

White dismantled both pillars of the majority's textual argument. First, the word ``sole'' appears in both the House and Senate Impeachment Clauses to divide prosecutorial from adjudicative authority between the two chambers of Congress---not to insulate the Senate from judicial review, just as Article I, \S~1's grant of ``All legislative powers'' does not preclude judicial review in cases like \textit{INS v. Chadha}, 462 U.S. 919 (1983) (\href{https://supreme.justia.com/cases/federal/us/506/224/}{506 U.S. at 241--42}). Second, ``try'' is clearly used in its legal sense, as confirmed by Article III, \S~2, cl.\ 3, and provides at least as manageable a standard as ``due process of law'' or ``Commerce \dots{} among the several States.''

Reaching the merits, however, White concluded that Senate Rule XI satisfies the Impeachment Trial Clause. Drawing on English parliamentary practice, Thomas Jefferson's \textit{Manual of Parliamentary Practice}, Justice Story's \textit{Commentaries}, federal courts' historical use of special masters, and the Senate's rulemaking authority under Article I, \S~5, cl.\ 2, White held that delegating initial factfinding to a Senate committee while reserving final argument and voting to the full Senate is fully compatible with the constitutional obligation to ``try'' impeachments (\href{https://supreme.justia.com/cases/federal/us/506/224/}{506 U.S. at 248--51}).

\subsection{Opinion Concurring in the Judgment (Justice Souter)}

Justice Souter agreed that Nixon's challenge was nonjusticiable, but grounded his conclusion in the prudential strands of \textit{Baker v. Carr} rather than a categorical textual bar (\href{https://supreme.justia.com/cases/federal/us/506/224/}{506 U.S. at 252--54}). Emphasizing ``the unusual need for unquestioning adherence to a political decision already made,'' the risk of ``embarrassment from multifarious pronouncements,'' and the severe governmental disruption judicial review would cause, Souter concluded that the Senate acted within the ``broad boundaries'' of its authority under Rule XI. Crucially, however, Souter reserved the possibility of judicial intervention if the Senate acted in a manner ``seriously threatening the integrity of its results''---such as convicting ``upon a coin toss, or upon a summary determination that an officer of the United States was simply `a bad guy'{}'' (\href{https://supreme.justia.com/cases/federal/us/506/224/}{506 U.S. at 253--54}).

\section{Examples: Future Applications}

\begin{enumerate}
\item \textbf{Same-Side (Nonjusticiable):} The Senate adopts an impeachment rule limiting defense cross-examination of witnesses to fifteen minutes per witness and excluding hearsay evidence through a presiding-officer ruling without a full-Senate vote. Under \textit{Nixon}, a removed official's lawsuit challenging these evidentiary and time limits as failing to constitute a proper ``trial'' is nonjusticiable, because Article I, \S~3, cl.\ 6 commits procedural trial design exclusively to the Senate and ``try'' supplies no judicially manageable evidentiary code.

\item \textbf{Same-Side (Nonjusticiable):} An impeached cabinet secretary challenges her conviction in federal court on the ground that the House of Representatives adopted articles of impeachment on a party-line vote without holding committee hearings or allowing her counsel to present exculpatory evidence. Because Article I, \S~2, cl.\ 5 grants the House the parallel ``sole Power of Impeachment,'' \textit{Nixon}'s textual commitment and separation-of-powers rationale bars judicial review of internal House impeachment procedures (\href{https://constitution.heritage.org/essays/a1-s2-c5-b/}{Heritage Guide}).

\item \textbf{Opposite-Side (Justiciable):} The House of Representatives votes to exclude a duly elected Representative who satisfies the age, citizenship, and residency requirements of Article I, \S~2, cl.\ 2, citing ethics allegations. As the \textit{Nixon} majority confirmed in preserving \textit{Powell v. McCormack}, 395 U.S. 486 (1969), this claim is justiciable because a separate constitutional provision establishes fixed, identifiable textual limits that courts can enforce without invading legislative discretion (\href{https://supreme.justia.com/cases/federal/us/506/224/}{506 U.S. at 237--38}).

\item \textbf{Opposite-Side (Justiciable):} The Senate votes 55--45 to convict and remove a federal appellate judge and declares the removal effective on the ground that a simple majority suffices for judicial impeachments. Unlike the open-ended word ``try,'' Article I, \S~3, cl.\ 6 expressly mandates ``the Concurrence of two thirds of the Members present.'' Under the majority's own distinction between imprecise grants of authority and ``identifiable textual limits'' (\href{https://supreme.justia.com/cases/federal/us/506/224/}{506 U.S. at 230, 238}), a conviction by less than two-thirds directly violates an unambiguous constitutional requirement and should be justiciable.

\item \textbf{Fence-Sitter (Borderline):} Following a House impeachment of a federal judge, the Senate skips all briefing, hearings, and committee proceedings and immediately votes 70--30 to convict based solely on a five-minute floor speech declaring the judge ``unfit.'' Under Rehnquist's categorical textual-commitment reasoning, even this summary conviction might remain nonjusticiable because ``try'' lacks judicially manageable standards and review threatens finality. Conversely, under Justice White's merits-review approach and Justice Souter's prudential exception for actions ``seriously threatening the integrity of its results'' (\href{https://supreme.justia.com/cases/federal/us/506/224/}{506 U.S. at 253--54}), federal courts could intervene because the Senate wholly abdicated any semblance of a trial.
\end{enumerate}

\section{Critique}

\textit{Nixon v. United States} stands as one of the rare modern Supreme Court decisions squarely invoking the political question doctrine to dismiss a constitutional claim (\href{https://www.everycrsreport.com/reports/R43834.html}{CRS Report R43834}; \href{https://review.law.stanford.edu/wp-content/uploads/sites/3/2023/05/Bradley-Posner-75-Stan.-L.-Rev.-1031.pdf}{Bradley \& Posner, 75 Stan.\ L.\ Rev.\ 1031 (2023)}). Scholarly reactions divide sharply along methodological and institutional lines.

\textbf{Originalist and Textualist Perspectives:} Many textualists and originalists praise Rehnquist's opinion for respecting the Constitution's explicit allocation of institutional authority. Commentators emphasize that ``sole'' appears only twice in the Constitution, that the Framers consciously rejected proposals to vest impeachment trials in the Supreme Court, and that structural separation of powers requires insulating Congress's only removal check on Article III judges from judicial self-protection (\href{https://constitution.heritage.org/essays/a1-s2-c5-b/}{Heritage Guide}). At the same time, other textualist-oriented scholars echo Justice White's critique that Rehnquist overread ``sole'' and underread ``try'': in context, ``sole'' allocated jurisdiction between the House and Senate rather than ousting judicial review, and declaring a constitutional verb wholly unenforceable risks transforming enumerated constitutional limits into parchment barriers (\href{https://www.law.cornell.edu/supct/html/91-740.ZC1.html}{White Concurrence}).

\textbf{Progressive and Doctrinal Perspectives:} Doctrinal and progressive scholars highlight the tension between \textit{Nixon}'s categorical abstention and \textit{Marbury v. Madison}'s baseline commitment to judicial enforcement of constitutional text. Scholars analyzing congressional regulation of the judiciary note that \textit{Nixon} confirms a domain of ``customary'' rather than judicially enforced constitutional norms, leaving impeachment procedures governed by political accountability, institutional tradition, and interbranch comity (\href{https://ilj.law.indiana.edu/articles/78/78_1_Geyh.pdf}{Geyh, 78 Ind.\ L.\ J.\ 153 (2003)}; \href{https://constitution.congress.gov/browse/essay/artII-S4-3/ALDE_00000698/}{Constitution Annotated}). Critics also observe that Justice Souter's prudential carve-out creates analytical instability: as Justice White noted, deciding whether the Senate stayed within ``broad boundaries'' already requires a preliminary merits judgment disguised as abstention (\href{https://scholarlycommons.law.hofstra.edu/cgi/viewcontent.cgi?article=2077&context=hlr}{Hofstra Law Review}).

\textbf{Competing Underlying Values:} At bottom, \textit{Nixon} reflects an enduring collision between \textit{rule-of-law accountability} (ensuring that no branch---not even the Senate---may disregard constitutional text or deprive an officer of procedural fairness) and \textit{separation-of-powers finality} (preserving democratic finality in high-stakes removals, preventing constitutional paralysis during a presidential impeachment, and avoiding a conflict of interest in which federal judges police the sole constitutional check on their own tenure).

\section{Key Quotations}

\begin{quote}
``The conclusion that the use of the word `try' in the first sentence of the Impeachment Trial Clause lacks sufficient precision to afford any judicially manageable standard of review of the Senate's actions is fortified by the existence of the three very specific requirements that the Constitution does impose on the Senate when trying impeachments: The Members must be under oath, a two-thirds vote is required to convict, and the Chief Justice presides when the President is tried.''
\end{quote}
\noindent --- Chief Justice Rehnquist, Opinion of the Court, 506 U.S. at 230 (\href{https://supreme.justia.com/cases/federal/us/506/224/}{Justia}; \href{https://www.courtlistener.com/opinion/112803/nixon-v-united-states/}{CourtListener}).

\begin{quote}
``Judicial involvement in impeachment proceedings, even if only for purposes of judicial review, is counterintuitive because it would eviscerate the `important constitutional check' placed on the Judiciary by the Framers.''
\end{quote}
\noindent --- Chief Justice Rehnquist, Opinion of the Court, 506 U.S. at 235 (\href{https://supreme.justia.com/cases/federal/us/506/224/}{Justia}; \href{https://www.law.cornell.edu/supct/html/91-740.ZO.html}{Cornell LII}).

\begin{quote}
``But we conclude, after exercising that delicate responsibility, that the word `try' in the Impeachment Trial Clause does not provide an identifiable textual limit on the authority which is committed to the Senate.''
\end{quote}
\noindent --- Chief Justice Rehnquist, Opinion of the Court, 506 U.S. at 238 (\href{https://supreme.justia.com/cases/federal/us/506/224/}{Justia}; \href{https://tile.loc.gov/storage-services/service/ll/usrep/usrep506/usrep506224/usrep506224.pdf}{U.S. Reports PDF}).

\begin{quote}
``Petitioner has not asked the Court to conduct his impeachment trial; he has asked instead that it determine whether his impeachment was tried by the Senate. The majority refuses to reach this determination out of a laudable desire to respect the authority of the Legislature. Regrettably, this concern is manifested in a manner that does needless violence to the Constitution.''
\end{quote}
\noindent --- Justice White, concurring in the judgment, 506 U.S. at 251 (\href{https://www.law.cornell.edu/supct/html/91-740.ZC1.html}{Cornell LII}; \href{https://supreme.justia.com/cases/federal/us/506/224/}{Justia}).

\begin{quote}
``If the Senate were to act in a manner seriously threatening the integrity of its results, convicting, say, upon a coin toss, or upon a summary determination that an officer of the United States was simply `a bad guy,' judicial interference might well be appropriate.''
\end{quote}
\noindent --- Justice Souter, concurring in the judgment, 506 U.S. at 253--54 (\href{https://supreme.justia.com/cases/federal/us/506/224/}{Justia}; \href{https://www.courtlistener.com/opinion/112803/nixon-v-united-states/}{CourtListener}).

\end{document}
